% Insert after the finite-horizon expression for Proposition 3.
\paragraph{Biased targets and measurement resolution.}
The same construction permits an imperfect candidate. Keep the observations fixed and multiply all three candidate revisions by $c$, with $0<c\le1$. Before the large region reaches the requested $cD$, its allowance follows the same recursion as in the perfect-candidate case. Once $cD$ is feasible, repeat releases of $cD$ replenish its allowance because $r^2D^2-(1-c)^2D^2>0$. The clipping of the two small corrections still cancels, so the large region and total both have per-issue absolute error $D[1-\min(c,s_t)]$. Their limiting RMSE is $(1-c)D$. The single-node allowance bound on the small downward revision still depends on the true $d_t$. Its actual downward release is at most $\kappa d_t$, so sharing the coefficient limits the large downward release to $\kappa D$. The third region can offset the total by at most $c\beta D$. The grouped regional and system lower bounds are therefore $(1-\kappa)D$ and $(1-\kappa-c\beta)D$. Both exceed the limiting full-projection score when $c(1-\beta)>\kappa$. This establishes robustness to proportional attenuation of the targets. If a change smaller than $m>0$ is unresolvable, all the small candidate corrections in the construction are visible only while $c\beta D\rho^{1-T}\ge m$. The horizon can then be at most $1+\lfloor\log(c\beta D/m)/\log\rho\rfloor$, provided $c\beta D\ge m$.

